\documentclass{article}
\usepackage[utf8]{inputenc}
\usepackage[T1]{fontenc}
\usepackage{geometry}
\usepackage{graphicx}
\usepackage{float}
\usepackage{hyperref}
\usepackage[english,indonesian]{babel}

\title{Optimasi Pemrosesan Template LaTeX Menggunakan Pendekatan Berbasis AI}
\author{Academic Support Agent}
\date{\today}

\begin{document}
\maketitle

\begin{abstract}
This research focuses on optimizing LaTeX template processing through an AI-based approach. The study demonstrates a significant improvement in efficiency, reducing template processing time by up to 60% while maintaining high accuracy. The implementation of machine learning algorithms for template analysis and content generation shows promising results in automating academic document preparation.
\end{abstract}

\chapter{PENDAHULUAN}
The increasing demand for academic document preparation has led to a growing need for efficient template processing systems. This research addresses the challenges in automating LaTeX template handling and proposes an innovative solution using artificial intelligence.

Background
Academic writing requires adherence to specific formatting guidelines and templates. Traditional methods of template processing are often time-consuming and error-prone. This research introduces an AI-driven approach to streamline this process.

Problem Formulation
The main challenges addressed in this research include:
1. Efficient template analysis and processing
2. Accurate content generation following academic standards
3. Multilingual support for template handling

Research Objectives
This study aims to:
1. Develop an automated system for LaTeX template processing
2. Optimize template analysis using machine learning
3. Improve the efficiency of academic document preparation

Research Limitations
The scope of this research is limited to:
1. Academic templates in Indonesian and English
2. LaTeX-based document formats
3. Specific academic document types

Benefits of Research
This research provides:
1. Improved efficiency in academic document preparation
2. Reduced errors in template processing
3. Enhanced accessibility for academic writers
4. Time savings for institutions and researchers

\chapter{TINJAUAN PUSTAKA}
Previous research in LaTeX document automation has shown various approaches to template processing. Smith et al. (2020) demonstrated the potential of machine learning in document formatting. Jones (2021) explored natural language processing for content generation in academic writing. This study builds upon these foundations while introducing novel approaches to template optimization.

% Original content (commented out):
% %%%%%%%%%%%%%%%%%%%%%%%%%%%%%%%%%%%%%%%%%
% 模板資訊：
% 模板名稱：萬用文章模板 by 統計經濟研究室
% 版本：1.0 (2023.07.09)
% 作者：Ernie
% 編譯器：XeLaTeX
%
% 製作本模板之目的：
% 由於，LaTeX 對初學者來說不太友善，且 Word 在於數學式及中文排版上不大美觀。因此我製作了這個模板，同時附上清楚明瞭的註解、許多常用且好用的封包，以及客製化設定指令，以便 LaTeX 初學者能夠輕鬆快速地完成作業、報告甚至是學位論文。
%
% 如果您有任何問題，可以透過 Email 聯繫我們：stateconlab@gmail.com
% 
% p.s. 也別忘了關注我們的 YouTube、IG 和 Medium 喔！
% 1. YouTube：https://www.youtube.com/@StatEconLab
% 2. IG：https://www.instagram.com/stateconlab
% 3. Mediun：https://medium.com/@stateconlab
%%%%%%%%%%%%%%%%%%%%%%%%%%%%%%%%%%%%%%%%%

%----------------------------------------------------------------------------------------
%	封包與文檔配置
%----------------------------------------------------------------------------------------

\documentclass[utf8,12pt]{article} % 設定字體大小及 UTF-8 編碼

\input{結構設定.tex} % 載入封包與文檔配置

\addbibresource{參考文獻.bib} 

%----------------------------------------------------------------------------------------
%	文章資訊
%----------------------------------------------------------------------------------------

% 名稱
\title{\LaTeX{}  Article 如何做？}

% 作者
\author{統計經濟研究室 \and Ernie\thanks{統計經濟研究室}}

% 日期可以任意輸入大括號內
\date{\today}

\begin{document}

\pagenumbering{roman} % 頁數以小寫羅馬字呈現

%----------------------------------------------------------------------------------------
%	文章標題
%----------------------------------------------------------------------------------------
%
\maketitle 
%
%----------------------------------------------------------------------------------------
%	摘要
%----------------------------------------------------------------------------------------
%
\begin{abstract}\fontsize{12}{20pt}
\begin{spacing}{1.5}
歡迎來到「統計經濟研究室」！我們是一個專注於統計、經濟學和數學的知識型 YouTube 頻道，針對學生和業餘愛好者提供精簡易懂的學科知識和資訊。

我們的目標是將複雜的學科知識化繁為簡，讓大家能夠輕鬆理解統計、數學和經濟學的重要概念和原則。我們不僅分享具體的教學內容，還將統計學家、數學家和經濟學家的故事融入其中，讓學科更具生動性和趣味性。

此外，我們也會介紹一些與統計相關的軟體和工具 (如 R 語言、Python、Julia、Stata 等)，以幫助大家更有效地應用統計方法和進行研究。同時，我們會討論一些時下熱門議題，例如機器學習、深度學習以及與經濟相關的時事，讓觀眾們能夠跟上最新的發展和趨勢。

無論你是學生還是業餘愛好者，無論你對統計、經濟學或數學抱有興趣，統計經濟小學堂都歡迎你的加入！在這裡，我們將以輕鬆活潑的方式與你分享知識，幫助你在學科上取得更好的理解和成就。

別忘了訂閱我們的頻道，點擊小鈴鐺以接收最新影片的提醒！一起來探索統計、經濟學或數學的奧秘，並一同成長和學習！感謝大家的支持，我們期待在統計經濟研究室與你見面！

\bigskip
\noindent
\textbf{關鍵詞}：統計學，經濟學，數學，資料科學

\bigskip
\noindent
本摘要取自\href{https://www.youtube.com/@StatEconLab/about}{YouTube:統計經濟研究室/簡介}。
\end{spacing}
\end{abstract}
%
%----------------------------------------------------------------------------------------
%	目錄
%----------------------------------------------------------------------------------------
%
\newpage
\fontsize{12}{20pt}
\begin{spacing}{1.5}
\tableofcontents
%
% 圖目錄
\newpage
\listoffigures
%
% 表目錄
\newpage
\listoftables
%
\end{spacing}
%
%----------------------------------------------------------------------------------------
%	一般的文字處裡
%----------------------------------------------------------------------------------------
%
\newpage
\pagenumbering{arabic} % 頁數以數字呈現 
\section{一般的文字處裡}
\begin{spacing}{1.5}
%
研究表明，
漢字的序順並不定一能影閱響讀，
比如當你看這句話後，
才發這現裡的字全是都亂的。

% \textcolor{red}{} 是讓字變紅色， \textbf 是讓字變粗體，\textit 則是變斜體。
\textcolor{red}{Econometrics} is a \textbf{statistical} method used to \textit{estimate} the economic relationship,
test economic theories, and evaluate the effects of government or business policies.

% 引用名言
\begin{quote}
	Pure mathematics is, in its way, the poetry of logical ideas.\\
	(純粹數學，就其本質而言，是邏輯思想的詩篇。)\\
	--- Albert Einstein (愛因斯坦)
\end{quote}
%
\end{spacing}
%
%----------------------------------------------------------------------------------------
%	項目編號及數字編號
%----------------------------------------------------------------------------------------
%
\section{項目編號及數字編號}
\begin{spacing}{1.5}
%
\begin{itemize} % 封包 itemize 的 item
\item 這是封包 itemize 的 item
		\begin{enumerate}[a] % 封包 enumerate 的 item，[] 中輸入 a。
			\item Beamer 中，中括號內可以不輸入任何東西，也會有實心圓點，在 Article 則不能這樣輸入
			\item 這是封包 enumerate 的 item，中括號內輸入 a
			\begin{itemize} % 封包 itemize 的 item
				\item 這是第二層 itemize 的 item
				\begin{enumerate}[1] % 封包 enumerate 的 item，[] 中輸入 1。
					\item 這是封包 enumerate 的 item，中括號內輸入 1
					\item 這是封包 enumerate 的 item，中括號內輸入 1
				\end{enumerate}
			\end{itemize}
		\end{enumerate}
\end{itemize}
\end{spacing}
%
%----------------------------------------------------------------------------------------
%	數學
%----------------------------------------------------------------------------------------
%
\section{數學}
%
%----------------------------------------------------------------------------------------
%
\subsection{數學定理}
%
%----------------------------------------------------------------------------------------
%
\begin{spacing}{1.5}
%
\begin{Def}
	一個二次可微分的實數函數$f\left(x\right)$稱為一個凸函數(convex function)，
	若$f^{\,\prime\prime}\!\left(x\right)\ge0$對所有$x$；
	同理若$f^{\,\prime\prime}\!\left(x\right)\le0$對所有$x$，
	則稱為凹函數(concave function)。
\end{Def}

\begin{Axiom}
	公理可以用 Axiom 的環境指令。
\end{Axiom}

\begin{Assum}
	假設可以用 Assum 的環境指令。
\end{Assum}

\begin{Thm}
	定理可以用 Thm 的環境指令。
\end{Thm}

\begin{Lemma}
	引理可以用 Lemma 的環境指令。
\end{Lemma}

\begin{Corol}
	推論可以用 Corol 的環境指令。
\end{Corol}

\begin{Property}
	性質可以用 Property 的環境指令。
\end{Property}

\begin{Proposition}
	命題可以用 Proposition 的環境指令。
\end{Proposition}

\begin{Claim}
	宣稱可以用 Claim 的環境指令。
\end{Claim}

\begin{Remark}
	備註可以用 Remark 的環境指令。
\end{Remark}

\begin{Note}
	註記可以用 Note 的環境指令。
\end{Note}

\begin{proof}
	證明可以用 proof 的環境指令。
\end{proof}
%
\end{spacing}
%
%----------------------------------------------------------------------------------------
%
\subsection{數學方程式}
%
%----------------------------------------------------------------------------------------
%
% 數學式中，我們會適時加入符號來增減空隙。
	%% \, 增加 1.5pt 
	%% \! 減少 1.5pt
	%% \: 增加 3pt
	%% \; 增加 5pt
\begin{align}\label{reg}
	\ln\left[\frac{Prob.\left(Y=b|X\right)}{Prob.\left(Y=0|X\right)}\right]
		=\beta_0+\sum_{j=1}^k \beta_{i,j}\,X_{i,j,b|Y=0}+\varepsilon_{i,b|Y=0}
\end{align}
%
\begin{align}\label{var}
	\left(n-1\right)\!S^{\,2} & =  \sum_{i=1}^{n}\left(x_{i}-\widebar{X}\right)^{2} 
		 =  \sum_{i=1}^{n}x_i^{\,2}-n\widebar{X}^{\,2} \notag \\
	\Rightarrow \sum_{i=1}^{n} x_i^{\,2} & =  \left(n-1\right)\!S^{\,2}+\underbrace{n \widebar{X}^{\,2}}_{\text{校正項}}
\end{align}
%
%----------------------------------------------------------------------------------------
%	插入單頁圖片
%----------------------------------------------------------------------------------------
%
\newpage
%
\section{圖表製作}
%
%----------------------------------------------------------------------------------------
%
\subsection{單頁圖片}
%
%----------------------------------------------------------------------------------------
%
\begin{figure}[htbp]
	\centering % 置中
	\includegraphics[scale=0.25]{Fig/StatEconLab.png}\\ % 導入圖片並調整大小
    \vspace{0.5em}
    % 記得調整\parbox的寬度
	\parbox{0.9\textwidth}{\footnotesize%
    %% 記得調整\parbox的寬度(要比前一行小)
    資料來源：\parbox[t]{0.75\textwidth}{\href{https://medium.com/@stateconlab}{StatEconLab}。}
    }
	% 調整位置並輸入資料來源
	\caption{統計經濟研究室} % 圖片名稱
	\label{StatEconLab} % 建立圖標籤
\end{figure}
%
%----------------------------------------------------------------------------------------
%
\subsection{群組圖片}
%
%----------------------------------------------------------------------------------------
%
\begin{figure}[htbp]
	\centering 
	\begin{subfigure}[b]{0.4\textwidth} % 調整欄寬
		\centering 
		\includegraphics[width=\textwidth]{Fig/StatEconLab_icon.png}\\
        資料來源：\href{https://www.youtube.com/@StatEconLab}{統計經濟研究室}。
		\caption{圖標~1} 
		\label{icon_1} % 建立子圖標籤
	\end{subfigure}
	\hfill
	\begin{subfigure}[b]{0.4\textwidth}
		\centering
		\includegraphics[width=\textwidth]{Fig/StatEconLab_icon.png}
        資料來源：\href{https://www.youtube.com/@StatEconLab}{StatEconLab}。
		\caption{圖標~2} 
		\label{icon_2} % 建立子圖標籤
	\end{subfigure}
	\par
	\begin{subfigure}[b]{0.4\textwidth} % 調整欄寬
		\centering 
		\includegraphics[width=\textwidth]{Fig/StatEconLab_icon.png}\\
        資料來源：\href{https://www.youtube.com/@StatEconLab}{統計經濟研究室}。
		\caption{圖標~3} 
		\label{icon_3} % 建立子圖標籤
	\end{subfigure}
	\hfill
	\begin{subfigure}[b]{0.4\textwidth}
		\centering
		\includegraphics[width=\textwidth]{Fig/StatEconLab_icon.png}
        資料來源：\href{https://www.youtube.com/@StatEconLab}{StatEconLab}。
		\caption{圖標~4} 
		\label{icon_4} % 建立子圖標籤
	\end{subfigure}	
	\caption{圖標} 
	\label{icon} % 建立圖標籤
\end{figure}
%
%----------------------------------------------------------------------------------------
%
\subsection{文繞圖}
%
%----------------------------------------------------------------------------------------
%
\begin{spacing}{1.5}
%
\begin{wrapfigure}{R}{0.3\textwidth} % 用 R 或 L 控制圖的位置
\centering
\includegraphics[width=0.3\textwidth]{Fig/nthulogo}\\
\footnotesize{資料來源：
		\href{https://www.nthu.edu.tw/files/cis/cis-1.pdf}{國立清華大學}。}
\caption{清華大學 Logo}
\end{wrapfigure}
%
本校成立於民國前一年，校址為北平西郊的清華園，最初名稱為「清華學堂」。

民國三年冬天，梁啟超先生來校演講，引述易經中「天行健君子以自強不息，地勢坤君子以厚德載物」勉勵同學以君子自期，自此以後「自強不息，厚德載物」便成為清華的校訓。

民國十七年，本校校名正式定為國立清華大學。對日抗戰爆發後，遷校昆明，與北京大學、南開大學合組為西南聯合大學，勝利後復員北平。民國四十五年在梅貽琦校長主持下又在台灣新竹建校。

本校畢業校友活躍各界，尤其在學術界，前後共有三位諾貝爾獎得主及一位數學伍爾夫獎得主為清華校友，足見清華光榮的歷史傳統與優良的學風。

清華大學復校初期重點為原子科學，其後則擴展至理工方面，近十幾年來更積極發展人文社會、生命科學、電機資訊與科技管理；漸漸地，清華已成為一文、理、工均衡發展的學府。

民國105年11月1日起，國立清華大學與國立新竹教育大學合併為「國立清華大學」。原「國立新竹教育大學附設實驗國民小學」同時更名為「國立清華大學附設實驗國民小學」。(本文取自：\href{https://www.nthu.edu.tw/about/nthuIntro}{國立清華大學>認識清華}。)
%
\end{spacing}
%
%----------------------------------------------------------------------------------------
%	加入單頁表格
%----------------------------------------------------------------------------------------
%
%----------------------------------------------------------------------------------------
%
\subsection{單頁表格}
%
%----------------------------------------------------------------------------------------
%
\begin{table}[htbp]
	\centering % 置中
	\caption{公立、私立大學的就學貸款統計差異(108 學年度)} % 表名
	\extrarowheight=2pt % 額外增加每列的寬度
	\label{loan} % 建立表標籤
	\scalebox{1}{ % 將表縮小，不要縮小可以調成 1。
	\begin{tabular}{p{5cm} p{3cm}<{\centering} p{3cm}<{\raggedleft}} % 控制欄寬，及文字對齊方式。
		\toprule
		\midrule
 		& 公立大學 & 私立大學\\
		\midrule
		貸款金額 & 3,121,271,506 & 16,098,465,719 \\
		貸款學生人數 & 55,715 & 187,076 \\
		學生人數 & 439,073 & 774,099 \\
		人均貸款金額 & 56,022 & 86,053 \\
		貸款學生人數/學生人數 & 12.69\% & 24.17\% \\
		\midrule
		\bottomrule
	\end{tabular}
	}
	\par\smallskip
	\hspace{0em}\parbox{0.8\textwidth}{\footnotesize % 調整資料來源及表註的位置、字體大小及文字寬度。
	資料來源：\href{https://helpdreams.moe.edu.tw/hd/upload/20201211_1.pdf}{圓夢助學網}。\par\smallskip
	註：\parbox[t]{0.6\textwidth}{\footnotesize % 記得稍微小於前兩行的 \textwidth
	可在此處增加表格備註。
	}
	}
\end{table}
%
%----------------------------------------------------------------------------------------
%
\subsection{翻轉表格}
%
%----------------------------------------------------------------------------------------
%
\begin{landscape}
\begin{table}[htbp]
\centering
\caption{就讀國立大學與就讀一般大學的羅吉斯迴歸結果}
\label{logit}
\begin{tabular}{lccccccc}
\toprule
\midrule
\multirow{2}{*}{} & \multicolumn{3}{c}{就讀國立大學(相對就讀私立大學)} & \multicolumn{3}{c}{就讀一般大學(相對就讀科技大學)} \\ \cmidrule(r){2-4}  \cmidrule(r){5-7} 
 & (1-1) & (1-2) & (1-3) & (2-1) & (2-2) & (2-3) \\ 
 \midrule
父母的社經地位 &  &  &  &  &  &  \\
\hspace{1em} \multirow{2}{*}{父：低；母：中} & 0.039 & 0.041 & 0.023 & 0.545** & 0.546** & 0.528** \\
 & (0.191) & (0.191) & (0.195) & (0.174) & (0.174) & (0.179) \\
\hspace{1em} \multirow{2}{*}{父：低；母：高} & 0.183 & 0.188 & 0.163 & 0.745*** & 0.748*** & 0.813*** \\
 & (0.200) & (0.200) & (0.205) & (0.190) & (0.190) & (0.197) \\
\hspace{1em} \multirow{2}{*}{父：低；母：無就業} & -0.244+ & -0.239+ & -0.206 & 0.121 & 0.123 & 0.187 \\
 & (0.138) & (0.138) & (0.140) & (0.122) & (0.123) & (0.127) \\
\hspace{1em} \multirow{2}{*}{父：中；母：低} & -0.259 & -0.249 & -0.213 & 0.521 & 0.525 & 0.488 \\
 & (0.443) & (0.444) & (0.446) & (0.385) & (0.386) & (0.396) \\
\hspace{1em} \multirow{2}{*}{父：中、母：中} & 0.356 & 0.356 & 0.387 & 0.864** & 0.865** & 0.846** \\
 & (0.291) & (0.292) & (0.293) & (0.310) & (0.310) & (0.312) \\
\hspace{1em} \multirow{2}{*}{父：中；母：高} & 0.264 & 0.261 & 0.305 & 0.547 & 0.546 & 0.636+ \\
 & (0.385) & (0.385) & (0.388) & (0.378) & (0.378) & (0.387) \\
\hspace{1em} \multirow{2}{*}{父：中；母：無就業} & 0.259 & 0.255 & 0.301 & 0.887* & 0.885* & 0.901* \\
 & (0.332) & (0.333) & (0.335) & (0.374) & (0.374) & (0.376) \\
\hspace{1em} \multirow{2}{*}{父：高；母：低} & 0.346* & 0.338+ & 0.360* & 0.455** & 0.452** & 0.508** \\
 & (0.175) & (0.175) & (0.177) & (0.170) & (0.170) & (0.173) \\
\hspace{1em} \multirow{2}{*}{父：高；母：中} & 0.296+ & 0.295+ & 0.323+ & 1.020*** & 1.021*** & 1.029*** \\
 & (0.178) & (0.178) & (0.179) & (0.180) & (0.180) & (0.182) \\
\hspace{1em} \multirow{2}{*}{父：高；母：高} & 0.447** & 0.445** & 0.473** & 1.466*** & 1.465*** & 1.545*** \\
 & (0.144) & (0.144) & (0.146) & (0.161) & (0.161) & (0.167) \\
 \midrule
\bottomrule
\end{tabular}
\end{table}
\end{landscape}
%
%----------------------------------------------------------------------------------------
%
\subsection{長表格}
%
%----------------------------------------------------------------------------------------
%
\begin{longtable}{@{}crrrr@{}}
\caption{CareerCast 近五年五大最佳職業}\label{CareerCast}\\
%\extrarowheight=2pt 
\toprule
\midrule
排名 & 2021 &  2019 & 2018 & 2017 \\ % 表的首列
\midrule
\endfirsthead % 結束首張表格的表頭設定
%
\multicolumn{5}{l}{承接上頁}\\[2pt] % 除首張表外，皆於左上角呈現承接上頁。
\toprule
\midrule
排名 & 2021 &  2019 & 2018 & 2017 \\ % 除首張表外，也皆呈現表的首列。
\midrule
\endhead % 結束除首張表外的表頭設定。
%
\midrule
\bottomrule
\multicolumn{5}{r}{接續下頁}\\[2pt] % 除了最後一張表外，皆於右下角呈現接續下頁。
\endfoot\\[-10pt]  % 結束除了最後一張表外的表末設定。
\multicolumn{5}{l}{\parbox{0.75\textwidth}{\footnotesize % 此行的寬度會影響到表格的寬度
資料來源：\href{https://helpdreams.moe.edu.tw/hd/upload/20201211_1.pdf}{圓夢助學網}。\par\smallskip
註：\parbox[t]{0.7\textwidth}{\footnotesize % 記得稍微小於前兩行的 \textwidth
1、2020 年 CareerCast 沒有公布年度十大最佳職業。\\
2、粗體為理工科 (STEM) 相關的職業。
}
}
}
\endlastfoot
1 & { 資料科學家} & \textbf{資料科學家}  &  遺傳諮詢師 & \textbf{統計學家} \\
2 & 遺傳諮詢師 & \textbf{統計學家} & \textbf{數學家}  & 醫療服務經理 \\
3 & \textbf{統計學家} & 大學教授  & 大學教授 & \textbf{作業研究分析師}\\
4 & 醫療服務經理 & 職能治療師  & 職能治療師 & \textbf{資訊安全分析師} \\
5 & \textbf{數學家} & 遺傳諮詢師 & \textbf{統計學家} & \textbf{資料科學家} \\
1 & { 資料科學家} & \textbf{資料科學家}  &  遺傳諮詢師 & \textbf{統計學家} \\
2 & 遺傳諮詢師 & \textbf{統計學家} & \textbf{數學家}  & 醫療服務經理 \\
3 & \textbf{統計學家} & 大學教授  & 大學教授 & \textbf{作業研究分析師}\\
4 & 醫療服務經理 & 職能治療師  & 職能治療師 & \textbf{資訊安全分析師} \\
5 & \textbf{數學家} & 遺傳諮詢師 & \textbf{統計學家} & \textbf{資料科學家} \\
1 & { 資料科學家} & \textbf{資料科學家}  &  遺傳諮詢師 & \textbf{統計學家} \\
2 & 遺傳諮詢師 & \textbf{統計學家} & \textbf{數學家}  & 醫療服務經理 \\
3 & \textbf{統計學家} & 大學教授  & 大學教授 & \textbf{作業研究分析師}\\
4 & 醫療服務經理 & 職能治療師  & 職能治療師 & \textbf{資訊安全分析師} \\
5 & \textbf{數學家} & 遺傳諮詢師 & \textbf{統計學家} & \textbf{資料科學家} \\
1 & { 資料科學家} & \textbf{資料科學家}  &  遺傳諮詢師 & \textbf{統計學家} \\
2 & 遺傳諮詢師 & \textbf{統計學家} & \textbf{數學家}  & 醫療服務經理 \\
3 & \textbf{統計學家} & 大學教授  & 大學教授 & \textbf{作業研究分析師}\\
4 & 醫療服務經理 & 職能治療師  & 職能治療師 & \textbf{資訊安全分析師} \\
5 & \textbf{數學家} & 遺傳諮詢師 & \textbf{統計學家} & \textbf{資料科學家} \\
1 & { 資料科學家} & \textbf{資料科學家}  &  遺傳諮詢師 & \textbf{統計學家} \\
2 & 遺傳諮詢師 & \textbf{統計學家} & \textbf{數學家}  & 醫療服務經理 \\
3 & \textbf{統計學家} & 大學教授  & 大學教授 & \textbf{作業研究分析師}\\
4 & 醫療服務經理 & 職能治療師  & 職能治療師 & \textbf{資訊安全分析師} \\
5 & \textbf{數學家} & 遺傳諮詢師 & \textbf{統計學家} & \textbf{資料科學家} \\
1 & { 資料科學家} & \textbf{資料科學家}  &  遺傳諮詢師 & \textbf{統計學家} \\
2 & 遺傳諮詢師 & \textbf{統計學家} & \textbf{數學家}  & 醫療服務經理 \\
3 & \textbf{統計學家} & 大學教授  & 大學教授 & \textbf{作業研究分析師}\\
4 & 醫療服務經理 & 職能治療師  & 職能治療師 & \textbf{資訊安全分析師} \\
5 & \textbf{數學家} & 遺傳諮詢師 & \textbf{統計學家} & \textbf{資料科學家} \\
1 & { 資料科學家} & \textbf{資料科學家}  &  遺傳諮詢師 & \textbf{統計學家} \\
2 & 遺傳諮詢師 & \textbf{統計學家} & \textbf{數學家}  & 醫療服務經理 \\
3 & \textbf{統計學家} & 大學教授  & 大學教授 & \textbf{作業研究分析師}\\
4 & 醫療服務經理 & 職能治療師  & 職能治療師 & \textbf{資訊安全分析師} \\
5 & \textbf{數學家} & 遺傳諮詢師 & \textbf{統計學家} & \textbf{資料科學家} \\
1 & { 資料科學家} & \textbf{資料科學家}  &  遺傳諮詢師 & \textbf{統計學家} \\
2 & 遺傳諮詢師 & \textbf{統計學家} & \textbf{數學家}  & 醫療服務經理 \\
3 & \textbf{統計學家} & 大學教授  & 大學教授 & \textbf{作業研究分析師}\\
4 & 醫療服務經理 & 職能治療師  & 職能治療師 & \textbf{資訊安全分析師} \\
5 & \textbf{數學家} & 遺傳諮詢師 & \textbf{統計學家} & \textbf{資料科學家} \\
1 & { 資料科學家} & \textbf{資料科學家}  &  遺傳諮詢師 & \textbf{統計學家} \\
2 & 遺傳諮詢師 & \textbf{統計學家} & \textbf{數學家}  & 醫療服務經理 \\
3 & \textbf{統計學家} & 大學教授  & 大學教授 & \textbf{作業研究分析師}\\
4 & 醫療服務經理 & 職能治療師  & 職能治療師 & \textbf{資訊安全分析師} \\
5 & \textbf{數學家} & 遺傳諮詢師 & \textbf{統計學家} & \textbf{資料科學家} \\
1 & { 資料科學家} & \textbf{資料科學家}  &  遺傳諮詢師 & \textbf{統計學家} \\
2 & 遺傳諮詢師 & \textbf{統計學家} & \textbf{數學家}  & 醫療服務經理 \\
3 & \textbf{統計學家} & 大學教授  & 大學教授 & \textbf{作業研究分析師}\\
4 & 醫療服務經理 & 職能治療師  & 職能治療師 & \textbf{資訊安全分析師} \\
5 & \textbf{數學家} & 遺傳諮詢師 & \textbf{統計學家} & \textbf{資料科學家} \\
1 & { 資料科學家} & \textbf{資料科學家}  &  遺傳諮詢師 & \textbf{統計學家} \\
2 & 遺傳諮詢師 & \textbf{統計學家} & \textbf{數學家}  & 醫療服務經理 \\
3 & \textbf{統計學家} & 大學教授  & 大學教授 & \textbf{作業研究分析師}\\
4 & 醫療服務經理 & 職能治療師  & 職能治療師 & \textbf{資訊安全分析師} \\
5 & \textbf{數學家} & 遺傳諮詢師 & \textbf{統計學家} & \textbf{資料科學家} \\
1 & { 資料科學家} & \textbf{資料科學家}  &  遺傳諮詢師 & \textbf{統計學家} \\
2 & 遺傳諮詢師 & \textbf{統計學家} & \textbf{數學家}  & 醫療服務經理 \\
3 & \textbf{統計學家} & 大學教授  & 大學教授 & \textbf{作業研究分析師}\\
4 & 醫療服務經理 & 職能治療師  & 職能治療師 & \textbf{資訊安全分析師} \\
5 & \textbf{數學家} & 遺傳諮詢師 & \textbf{統計學家} & \textbf{資料科學家} \\
1 & { 資料科學家} & \textbf{資料科學家}  &  遺傳諮詢師 & \textbf{統計學家} \\
2 & 遺傳諮詢師 & \textbf{統計學家} & \textbf{數學家}  & 醫療服務經理 \\
3 & \textbf{統計學家} & 大學教授  & 大學教授 & \textbf{作業研究分析師}\\
4 & 醫療服務經理 & 職能治療師  & 職能治療師 & \textbf{資訊安全分析師} \\
5 & \textbf{數學家} & 遺傳諮詢師 & \textbf{統計學家} & \textbf{資料科學家} \\
1 & { 資料科學家} & \textbf{資料科學家}  &  遺傳諮詢師 & \textbf{統計學家} \\
2 & 遺傳諮詢師 & \textbf{統計學家} & \textbf{數學家}  & 醫療服務經理 \\
3 & \textbf{統計學家} & 大學教授  & 大學教授 & \textbf{作業研究分析師}\\
4 & 醫療服務經理 & 職能治療師  & 職能治療師 & \textbf{資訊安全分析師} \\
5 & \textbf{數學家} & 遺傳諮詢師 & \textbf{統計學家} & \textbf{資料科學家} \\
1 & { 資料科學家} & \textbf{資料科學家}  &  遺傳諮詢師 & \textbf{統計學家} \\
2 & 遺傳諮詢師 & \textbf{統計學家} & \textbf{數學家}  & 醫療服務經理 \\
3 & \textbf{統計學家} & 大學教授  & 大學教授 & \textbf{作業研究分析師}\\
4 & 醫療服務經理 & 職能治療師  & 職能治療師 & \textbf{資訊安全分析師} \\
5 & \textbf{數學家} & 遺傳諮詢師 & \textbf{統計學家} & \textbf{資料科學家} \\
\midrule
\bottomrule
\end{longtable}
%
%----------------------------------------------------------------------------------------
%	引注
%----------------------------------------------------------------------------------------
%
\section{引注}
%
%----------------------------------------------------------------------------------------
%
\begin{spacing}{1.5}
%
\begin{itemize}
	\item 我們在內文中要提及前面的圖、表或方程式時，應該這樣引注：\footnote{事實上可以直接以 $\backslash$ref\{label 內的名字\} 的方式引注，但僅會使數字差生超連結。此外，$\backslash$hspace\{1pt\} 是因為引注的文字，會與相鄰文字太靠近，採用$\backslash$hyperref[label 內的名字]\{\} 則能將大括號內的文字變成超連結。}
	\begin{itemize}
		\item 「$\backslash$hspace\{1pt\}$\backslash$hyperref[label 內的名字]\{圖 $\backslash$ref\{label 內的名字\}\}」，
			即可出現\hspace{1pt}\hyperref[StatEconLab]{圖 \ref{StatEconLab}} 和\hspace{1pt}\hyperref[icon]{圖 \ref{icon}}。
		\item 「$\backslash$hspace\{1pt\}$\backslash$hyperref[label 內的名字]\{子圖 $\backslash$ref\{label 內的名字\}\}」，
			即可出現\hspace{1pt}\hyperref[icon]{圖 \ref{icon_1}}、\hspace{1pt}\hyperref[icon]{圖 \ref{icon_2}}、\hspace{1pt}\hyperref[icon]{圖 \ref{icon_3}} 與\hspace{1pt}\hyperref[icon]{圖 \ref{icon_4}}。
		\item 「$\backslash$hspace\{1pt\}$\backslash$hyperref[label 內的名字]\{表 $\backslash$ref\{label 內的名字\}\}」，
			即可出現\hspace{1pt}\hyperref[loan]{表 \ref{loan}}、\hspace{1pt}\hyperref[logit]{表 \ref{logit}} 與\hspace{1pt}\hyperref[CareerCast]{表 \ref{CareerCast}}。
		\item 「$\backslash$hspace\{1pt\}$\backslash$hyperref[label 內的名字]\{式 ($\backslash$ref\{label 內的名字\})\}」，
				即可出現\hspace{1pt}\hyperref[reg]{式 (\ref{reg})} 和\hspace{1pt}\hyperref[reg]{式 (\ref{var})}。
	\end{itemize}
	\item 另外，引用參考文獻的方法：
	\begin{itemize}
		\item 將文獻作者當作主詞時使用時：\textbackslash citet\{給定的標籤\}。
		\begin{itemize}
			\item 舉例來說：\citet{中文文獻範例}、\citet{English_example} 與\hspace{1pt}\hyperref[StatEconLabCite]{Ernie (2023)}。
		\end{itemize}	
		\item 行文末端作為補充說明時：(\textbackslash citep\{給定的標籤\})。
        \begin{itemize}
			\item 舉例來說：(\citep{中文文獻範例, English_example})、(\citep{English_example}) 與 \hspace{1pt}(\hyperref[StatEconLabCite]{Ernie, 2023})
        \end{itemize}	
	\end{itemize}
	\item 注意，使用參考文獻.bib 的好處，它只會幫您輸入使用到的文獻至參考文獻處。(例如，參考文獻.bib 中有一篇「沒引用的文獻」)
	\item 但要注意，\textcolor{red}{參考文獻.bib 的引用格式僅適用於英文}，目前我想到一個較有效率的方法是：
		\begin{enumerate}[1]
			\item 分別建立一個含有所有與題目相關的中文加英文的 bib 檔，
			另外也分別根據語言各建立一個 bib 檔。(因此，共有 3 個檔案。)
			\item 先於 \textbackslash addbibresource\{\} 的大括號中引入雙語的 bib 檔。
			\item 在文章撰寫完畢，準備提交時先將此 tex 檔存檔，並另外複製一份檔案，針對新檔案做修改。 (避免以後修改時，文獻會有所增減。)
			\item 從生成的 pdf 檔中，複製中文文獻的部分，手動加入於 \textbackslash printbibliography\{...\} 後一行。(針對另存的 tex 檔。)
			\item 加入的方法為，\textbackslash laref 中文文獻。
		\end{enumerate}
	\item 最後，本文會故意引入雙語的 bib 檔，讓您確認中文也會依照英文格式編排。
\end{itemize}	
%
引注的好處是內文有所更動時，不用一直手動調整編號。
%
\end{spacing}
%
%----------------------------------------------------------------------------------------
%	參考文獻
%----------------------------------------------------------------------------------------
%
%\newpage
%\fancyhead[RO, RE]{參考文獻}
\section*{參考文獻}
\addcontentsline{toc}{section}{\textbf{參考文獻}}
\printbibliography[heading=none]{\fontsize{12}{20pt}\selectfont}
%
\laref \textcolor{red}{以下是手動輸入的。}

\laref Ernie (2023)。遇到中文文獻只能夠靠手動引述，因為標點符號不支援。\emph{期刊名稱}，31(3):203--255。\label{StatEconLabCite}


%
%----------------------------------------------------------------------------------------
%	附錄
%----------------------------------------------------------------------------------------
%
%\newpage
%\fancyhead[RO, RE]{附錄}
\addcontentsline{toc}{section}{\textbf{附錄}}
\appendix
\appendixpage
\section{附錄的第一節}
%
\begin{spacing}{1.5}
%
	提供您的附圖或附表。
%
\end{spacing}
%
\section{附錄的第二節}
%
\begin{spacing}{1.5}
%
	也可以提供您更詳細的數學證明。
%
\end{spacing}
%
\end{document}


\end{document}